\section{相关工作}

\subsection{LLM 驱动的数据分析系统}

现有 LLM 数据分析系统大多将重点放在执行阶段，包括任务拆解、代码生成、执行纠错和结果解释等能力。DS-Agent、Data Interpreter、Data-to-Dashboard 等系统已经证明 LLM 可以显著提升数据分析自动化程度\citep{guo2024dsagent,hong2025data,zhang2025data}；面向数据科学代理的综述也表明，这一方向正在从“代码助手”向“分析代理”演进\citep{sun2025survey,zhao2024llm}。但这类方法在“给定问题之后如何把分析做出来”这一层面取得进展的同时，对规划阶段的关注仍相对不足：规划起点常常是自然语言问题本身，而不是当前数据集的真实结构与证据分布，因此在复杂结构化任务或跨领域任务上容易出现规划与数据事实脱节的问题。就本文关心的结构化分析规划而言，单纯增强执行器并不能自动保证规划质量，尤其不能保证计划具有足够的变量锚点、方法组织与证据闭环。

\subsection{多智能体协作与规划反馈}

多智能体框架为复杂分析任务提供了模块化协作路径。AutoGen、MetaGPT、CAMEL 等框架说明，通过把规划、规格化、执行和评审拆分到不同角色，系统可以获得更清晰的职责边界和更稳定的错误定位能力\citep{wu2023autogen,hong2024metagpt,li2023camel}。与此同时，ReAct 代表了单智能体“思考-行动-观察”循环的典型实现\citep{yao2023react}，Reflexion 则代表了基于执行反馈的再尝试机制\citep{shinn2024reflexion}。然而，已有工作常常仍以执行正确性为中心，对“首轮执行之后如何把反馈重新写回规划”讨论不足。本文关注的两轮迭代机制，正是试图把执行反馈从“纠错信号”提升为“再规划信号”。因此，本文实验中的 \texttt{G\_react\_single\_agent} 和 \texttt{H\_execution\_feedback} 被明确定位为 literature-inspired external mechanism baselines：前者对标 ReAct 风格单智能体规划族，后者对标 Reflexion / Self-Refine 风格的执行反馈机制族，而不是对某一篇论文系统的逐字复现。

\subsection{知识增强与图谱约束}

知识图谱增强大语言模型已被广泛用于检索、问答和内容生成\citep{liu2016knowledge,pan2024unifying,edge2024graphrag,mialon2023augmented,gao2023retrieval}，但在分析规划任务中，知识作用方式并不应停留在“提供补充事实”层面。对于结构化分析任务，系统需要区分两类不同约束：一类约束哪些变量关系值得优先考察，另一类约束在特定证据条件下哪些分析方法更可执行。本文将二者分别映射为因果图谱与方法图谱，目的不是降低模型自由度本身，而是在自由生成与结构约束之间建立更可控的中间层。与只做检索增强的路线不同，本文强调知识资源在规划阶段的作用是“约束决策结构”，而不是简单拼接到上下文中。

\subsection{本文定位}

与上述工作相比，本文的差异不在于单独引入某一项能力，而在于把“数据证据、知识约束和执行反馈”统一组织为一个规划框架。本文关注的是：系统如何在执行之前形成更好的分析方案，以及为什么这种方案在重复实验、外部机制基线比较、子领域稳定性验证和跨领域小规模验证下能够保持更高的稳定性。换言之，本文定位不是“再实现一个能生成代码的分析系统”，而是提出一个针对结构化分析规划阶段的可复验软件方法框架；专利分析是主要验证场景，软件工程 Bugzilla 数据则用于提供额外的外部有效性信号。
